\intro

Численное моделирование теплопередачи широко применяется при проектировании изделий машиностроения, энергетики и электроники, где температурный режим определяет надёжность конструкции~\cite{krainov2016}. Описание процесса строится на уравнении теплопроводности Фурье--Кирхгофа~\cite{chichko2013mathematical}, которое для тел сложной формы и при неоднородных граничных условиях, как правило, не имеет аналитического решения. Поэтому на практике используют сеточные методы, в которых расчётная область заменяется множеством узлов, а производные --- конечными разностями~\cite{kolgan1975finite}.

В предыдущем семестре был разработан программный комплекс HeatSim3D, который решает трёхмерную нестационарную задачу теплопроводности по явной конечно-разностной схеме FTCS. Комплекс состоит из расчётного модуля на языке C++~\cite{stroustrup1986overview} и графического интерфейса на языке Python~\cite{python2021python}. Работа с ним показала, что время расчёта быстро растёт при измельчении сетки. При уменьшении шага сетки вдвое число узлов увеличивается в восемь раз, а допустимый шаг по времени из условия устойчивости уменьшается в четыре раза, так что общий объём вычислений возрастает примерно в тридцать два раза. Для сетки из нескольких миллионов узлов один расчёт в последовательном режиме занимает уже минуты, а серия расчётов для разных материалов и граничных условий --- часы.

Современные персональные компьютеры оснащены многоядерными процессорами, однако последовательная программа использует только одно ядро. Наиболее простой способ задействовать все ядра в программе на C++ даёт технология OpenMP~\cite{dagum1998openmp}, которая позволяет распараллелить существующий код добавлением директив компилятора без изменения структуры алгоритма~\cite{chandra2001parallel}. При этом важно не только получить выигрыш во времени, но и доказать, что параллельная программа выдаёт те же результаты, что и последовательная, а также оценить, насколько полно используются вычислительные ресурсы. Именно количественная оценка ускорения и анализ факторов, которые его ограничивают, составляют исследовательскую часть работы.

Целью работы является разработка и исследование параллельной OpenMP-версии программного комплекса HeatSim3D для численного моделирования нестационарной теплопередачи в сплошных средах.

Для достижения поставленной цели необходимо решить следующие задачи:
\begin{enumerate}[label=\arabic*.]
    \item Изучить технологию OpenMP и показатели эффективности параллельных программ для систем с общей памятью.
    \item Проанализировать информационные зависимости алгоритма явной схемы FTCS и выделить участки, допускающие параллельное выполнение.
    \item Разработать параллельную версию расчётного модуля с управляемым числом потоков, совместимую с компиляторами GCC, MinGW и Microsoft Visual C++.
    \item Обеспечить сборку и работу программного комплекса в операционных системах Linux и Windows.
    \item Проверить корректность параллельной версии путём сравнения с аналитическим решением и с результатами последовательного расчёта.
    \item Исследовать ускорение и эффективность параллельной версии для расчётных сеток разного размера и оценить влияние последовательной записи результатов.
\end{enumerate}

Новизна работы для программного комплекса HeatSim3D состоит в том, что параллельный режим впервые снабжён средствами управления числом потоков и автоматического замера производительности. Это позволило перейти от качественного утверждения об ускорении расчёта к измеренным значениям ускорения и эффективности. Кроме того, разработан модуль построения графиков в оформлении, пригодном для научного отчёта, и набор автоматических тестов, подтверждающих корректность параллельной версии.

В первой главе приводится краткое описание математической модели и численного метода, рассматриваются модель исполнения OpenMP и показатели эффективности параллельных программ. Во второй главе описана разработка параллельной версии расчётного модуля, включая анализ зависимостей по данным, выбор директив OpenMP и обеспечение кроссплатформенности. В третьей главе представлены результаты проверки корректности и исследования производительности параллельной версии на многоядерном процессоре.
